\def\xhref#1#2{\href{#1}{\hskip16pt}\llap{#2}}
\def\xxhref#1#2{\href{#1}{\hskip14pt}\llap{#2}}
\makeatletter
\def\@oddhead{\hfill}
\def\@evenhead{\hfill}
\makeatother
   
   
\normalsize
%\fontsize{9.6}{11.6}\selectfont
%\parskip4pt plus2pt minus1pt

%\advance\baselineskip by-.1pt

\pznspis{Nieunikniona kolizja?}


\pzntyt{Nieunikniona kolizja?}

\vspace*{.6\parskip}%
\begin{dwieszpalty}\advance\baselineskip by-.2pt
  Przez wiele lat astronomowie uważali, że Droga Mleczna
zderzy się ze swoją galaktyką sąsiadką
%zderzy się z~najbliższą sąsiadką, dużą galaktyką
--  Andromedą. Symulacje na podstawie nowych danych wykazują jednak, że prawdopodobieństwo wystąpienia tej kolizji to tylko 50$\%$, przynajmniej w~ciągu następnych 10~miliardów~lat. Co się zmieniło?

Nasza galaktyczna ,,dzielnica'' Wszechświata zwana jest Grupą Lokalną. Znajdują się w~niej dwie duże galaktyki spiralne: nasza Droga Mleczna i~galaktyka Andromedy (zwana też M31) oraz około 100 znanych mniejszych galaktyk. \marg{W~Grupie Lokalnej jest prawdopodobnie wiele małych galaktyk, które nie zostały jeszcze odkryte.} Dotychczasowe badania Grupy Lokalnej sugerowały, że nasza Galaktyka i~Andromeda są na kursie kolizyjnym. Zderzenie dwóch galaktyk miałoby nastąpić za około 4 miliardy lat, a~kolejne zderzenia w~odstępach 2~miliardów~lat dałyby początek nowej, dużej galaktyce eliptycznej (powstałej z~połączenia Andromedy i~Drogi Mlecznej) pełnej gwiazd rozproszonych na nowych orbitach. Ziemia, Słońce i~pozostałe planety naszego Układu Słonecznego przetrwałyby zderzenie, ale cały Układ Słoneczny znalazłby się na obrzeżach nowo powstałej galaktyki eliptycznej. 

Te badania mają jednak już kilka dobrych dekad. W~międzyczasie poszerzyła się nasza wiedza na temat mas, orbit oraz dokładnych prędkości galaktyk należących do Grupy Lokalnej. Wykonaliśmy też nowe obserwacje z~wykorzystaniem europejskiej sondy kosmicznej Gaia oraz Kosmicznego Teleskopu Hubble'a. Naukowcy na podstawie nowych badań %i obserwacji
zoptymalizowali i~poprawili symulacje szacujące prawdopodobieństwo zderzenia Drogi Mlecznej z~Andromedą. Wzięli też pod uwagę, wcześniej zaniedbywany, możliwy wpływ innych galaktyk Grupy Lokalnej na kurs kolizyjny Andromedy i~Drogi Mlecznej. Okazało się, że był to strzał w~dziesiątkę. Odkryto bowiem, że w~szczególności dwie pobliskie galaktyki -- satelity (M33 i~Wielka Mgławica Magellana) mogą mieć kolosalny wpływ na to, czy do zderzenia w~ogóle dojdzie! Uwzględnienie tylko tych dwóch galaktyk zmniejszyło prawdopodobieństwo zderzenia Drogi Mlecznej i~Andromedy w~ciągu najbliższych 10~miliardów lat do blisko 50$\%$ -- w~połowie scenariuszy do zderzenia nie doszło.

Oczywiście autorzy publikacji prezentującej wyniki badań są ostrożni i~twierdzą, że modele i~symulacje muszą być jeszcze udoskonalone, aby można było dokładnie ustalić, jakie będą trajektorie ruchu galaktyk, a~w~konsekwencji, z~jakim prawdopodobieństwem Andromeda zderzy się z~naszą Galaktyką. Mogą nam pomóc nadchodzące obserwacje z~sondy Gaia, które dostarczą informacji o~ruchach właściwych i~masie galaktyk Grupy Lokalnej. Okazało się, że nieunikniona katastrofa wcale taka nieunikniona nie jest.
\end{dwieszpalty}

\vskip-2pt

\marg[22pt]{Możliwe orbity ruchu Drogi Mlecznej (MW) i~Andromedy (M31). Kolorowe linie pokazują różne możliwe trajektorie ruchu galaktyk w~ciągu najbliższych 10~miliardów lat lub do momentu zderzenia, jeżeli nastąpi ono wcześniej. Na rysunkach po lewej stronie trajektorie ruchu bez uwzględnienia wpływu innych galaktyk Grupy Lokalnej, po prawej z~uwzględnieniem wpływu galaktyk M33 i~Wielkiego Obłoku Magellana. Wykresy w~górnym panelu pokazują widok ,,z~góry'', a~wykresy na dole widok ,,od boku''. Rysunek pochodzi z~arXiv:2408.00064.
  \vskip.8cm

{Artykuł oparty na publikacji: Till Sawala, Jehanne Delhomelle, Alis J. Deason et al., ,,Apocalypse When? No Certainty of a~Milky Way -- Andromeda Collision'',  arXiv:2408.00064. }}\includegraphics[scale=0.25]{z–góry.png}\\[-6pt]
\includegraphics[scale=0.25]{z-boku.png}
\medskip
% Źródła
% https://www.pnas.org/doi/10.1073/pnas.2409983121
% https://www.bbc.com/news/articles/czxl849j77ko





\rightline{\large\textit{Anna DURKALEC}%
%\textit{Michał BEJGER}%
%\textit{Katarzyna MAŁEK}%
}%
%}

\medskip

%{\scriptsize \rightline
%  {Centrum Astronomiczne im. Mikołaja Kopernika PAN,}
%   \rightline  {Istituto Nazionale di Fisica Nucleare (INFN), Sezione di Ferrara, Włochy}\par} 

{\scriptsize \rightline{Zakład Astrofizyki, Departament Badań Podstawowych, Narodowe Centrum Badań Jądrowych}}  %A.Durkalec

 %\rightline{ \scriptsize  Departament Badań Podstawowych (BP4), Zakład Astrofizyki, Narodowe Centrum Badań Jądrowych}%K.Małek
%\vfill
% \end{dwieszpalty}
\vfill

\eject

\normalsize
%\input 11-niebo




\endinput